\documentclass[a4paper]{arrowhead}
\usepackage[T1]{fontenc}
\usepackage{longtable,booktabs,xurl,enumitem}
\setlength{\emergencystretch}{3em}
\setlength{\parindent}{0pt}\setlength{\parskip}{5pt}
\setlist{nosep,leftmargin=*}
\usepackage{xurl}
\begin{document}
\ArrowheadTitle{LocalCloudManager / SysD (Go)}\ArrowheadType{System Description}\ArrowheadTypeShort{Go}
\ArrowheadVersion{5.3.0}\ArrowheadDate{2026-11-09}
\ArrowheadAuthor{Eclipse Arrowhead migration project}\ArrowheadStatus{REVIEW DRAFT}
\ArrowheadContact{Project maintainers}\ArrowheadFooter{Eclipse Arrowhead -- Go LocalCloudManager -- EPL-2.0}
\ArrowheadSetup
\section{Documentation edition and implementation scope}
This is the 5.3.0 documentation edition, dated 2026-11-09. It carries forward the system and service contracts from the retained 5.2.2 documentation and identifies the current primary implementation sources. The former \path|documentation/v5.2.2| folder is retained unchanged.

Document version 5.3.0 is distinct from executable, dependency, registered system/service and wire-contract versions. Values still defined by the source (including 5.2.2 artefact versions, 1.0.0 service versions and \path|SysMonitor-AMQP-IDD-5.2.2| payload identifiers) are preserved. This documentation update does not change runtime behaviour or establish a new implementation release. Historical verification statements describe the tests recorded in the inherited baseline; this revision validates document compilation and references.

\section{System overview}
LocalCloudManager is the operator-facing Go adapter for an Arrowhead local cloud. This package relocates the latest located local source into \path|ah5-support-go/LocalCloudManager|, with a dedicated launcher, bundled shared-runtime dependencies and embedded dashboard assets. The target system version is 5.2.2; produced service versions are 1.0.0. This relocation does not establish complete Java interchangeability or modify a deployed cloud.
\section{System boundary and stakeholders}
An operator browser uses the console UI and localCloudManagement API. The adapter consumes ServiceRegistry, Authentication, Simple Store Service Orchestration, DataManager, SysMonitor and PlantDescriptionEngine. It reads Docker container metadata/logs and performs limited lifecycle actions through a Unix socket. Native Go business code does not invoke a Java, Rust or Node.js server.
\section{Produced services}
\path|localCloudManagement| at \path|/api| is the browser-facing session API. Historical \path|monitor| at \path|/localcloud/monitor| provides echo and retains the source's legacy registration name. Optional \path|sysMonitor| is native AMQPS telemetry and is enabled only in secure mode with \path|SYSMONITOR_ENABLED=true|. Legacy monitor naming is not a recommendation for new telemetry contracts.
\section{Consumed services and prerequisites}
Provision a distinct LocalCloudManager sy certificate, matching private key, trusted CA bundle, Authentication identity/password, and provider grants/management whitelist. Configure explicit HTTPS URLs for every used provider: several inherited defaults are HTTP and fail in secure mode. Docker access requires the intended Compose project and isolated socket access. No CA private key is needed or included.
\section{Deployment and variants}
Exactly one manager implementation should serve a cloud. The revised SysML model records Go's canonical source, the existing Java-cloud Node.js adapter, and the Rust source as distinct variants; it does not assert tested substitution. Default standalone port8480 uses mutual TLS and serves UI/API on the same listener. Historical API9480/UI9490 mappings belong to deployment configuration, not this relocation. Existing native-go deployment source is retained unchanged to preserve shared-runtime build contexts.
\section{Compatibility and release gates}
This relocation intentionally preserves the latest located manager's business behavior. New changes are packaging (standalone nested module/role-specific launcher), embedded UI fallback, documentation and tests. This is not a security remediation or full Java-equivalence implementation.\par\medskip
1. \path|PUT/DELETE /api/sysmonitor/subscriptions/{providerSystem}| currently checks confirmation and returns updated/removed flags without forwarding a write. Do not interpret this as applied configuration. This remains a release blocker.\par\medskip
2. \path|/api/intercloud| returns a static unavailable/no-relay summary. It does not inspect Gatekeeper/GateTunnel. Consumer cloud labels include hard-coded TestCloud strings; no live foreign authorization catalogue is queried.\par\medskip
3. Overview hard-codes registration/naming flags and leaves several collections empty. Subscription status hard-codes intake.connected=true/providerFound=false. These are display adapters, not authoritative readiness or security evidence.\par\medskip
4. Historian, SysMonitor, plant and orchestration adapters target the inherited Go-cloud routes/shapes. They are not proven compatible with the newer generic Go support port or latest Java APIs. The standalone module must be qualified against the providers actually deployed.\par\medskip
5. Console Host/Origin checks compare host names, not exact origins/ports. The \path|ah_console| cookie is not port-isolated and lacks Secure. Sessions/history are in-memory and reset on restart. Missing password disables APIs. Login lockout is global, five failures in five minutes then one-minute lockout; this is not comprehensive abuse protection.\par\medskip
6. The manager API is dispatched before the shared application's token/management checks. Its boundary is console password/session plus the transport's mTLS if enabled, not a bearer-token-only API. Static assets and health have no console session requirement. A proxy must preserve this intended boundary.\par\medskip
7. Authentication has explicit CN/profile pinning. Other provider clients use CA/SAN validation but do not generally pin expected provider CN. Do not claim universal authority-substitution resistance for this inherited code. Redirects are disabled. Optional CRLs reload for new connections only; omitted CRLs mean no revocation enforcement.\par\medskip
8. Lifecycle uses the static systemToCompose service allowlist and Compose-project filtered container list; MANAGED\_SERVICES is not consumed by this implementation. Action is start/stop/restart only. Logs, however, can read any returned project container, not just the lifecycle allowlist. The raw Docker socket still confers broad daemon privilege: application checks are not a least-privilege socket boundary. Production needs a constrained proxy and deployment review.\par\medskip
9. Shared runtime source includes other roles for compile-time dependencies and inherited tests. The dedicated executable refuses selecting those roles; they are not a second support-system deployment. No Docker socket or live state is needed by tests.\par\medskip
10. Telemetry AMQP publisher has no durable local outbox; retries/reconnects do not establish exactly-once delivery. power=line/available and process health are generated status, not measurements of physical hardware. Registration is retried every20seconds; displayed registration flags are not confirmation. No real broker test is claimed here.\par\medskip
11. Shared state is the inherited JSON format, not the newer support port's transactional state model. There is no new Java SQL migration, encryption-at-rest or multi-process ownership guarantee. Keep independent state and backups.\par\medskip
12. Historical Markdown/SysML claims of three-way interchangeability are retained only under docs/legacy. The revised model and new SysD/SysDD/SD/IDD set distinguish intended common contracts from actual Go qualification.\par\medskip
13. The shared serviceOperations helper generates POST/name metadata for localCloudManagement and the legacy monitor service, whereas actual handlers include GET/PUT/PATCH/DELETE and nested paths. Registration metadata is therefore not an authoritative HTTP contract. The new IDDs document actual handler methods/paths; correcting advertised metadata is separate functional work, not silently folded into relocation.\par\medskip
Source relocation does not modify running containers or remove their deployment source. Before switching deployment build paths, close these gaps, qualify provider DTOs/policies, review Docker access, isolate browser trust domains, and test rollback.\par\medskip
\section{References and verification}
Authoritative moved code: LocalCloudManager/internal/cloud and embedded ui sources. Exact original/relocated SHA256 hashes are in docs/source-manifest.json. docs/legacy preserves the earlier local Markdown/SysML, not an endorsement of its parity claims. The revised SysML variability model is descriptive and has not been checked by a SysML parser. Executed checks and qualification limits are recorded in VERIFICATION.md. Source movement does not deploy or alter running clouds.
\section{Revision record}
2026-11-09: Documentation edition 5.3.0 derived from the retained 5.2.2 sources; metadata and local references updated. Runtime and wire-contract versions remain source-defined.
\end{document}
